\documentclass[11pt]{article}
\usepackage[margin=25mm]{geometry}
\usepackage{fontspec}
\setmainfont{Latin Modern Roman}
\usepackage{amsmath,amssymb,mathtools}
\usepackage{unicode-math}
\setmathfont{Latin Modern Math}
\usepackage{xurl,hyperref}
\hypersetup{colorlinks=true,urlcolor=blue}
\setcounter{secnumdepth}{0}
\setlength{\parindent}{0pt}
\setlength{\parskip}{5pt}
\setlength{\emergencystretch}{3em}
\newcommand{\sref}[2]{\hyperlink{source:#1:#2}{[S#2]}}
\newcommand{\cataloguescope}[1]{\paragraph{Recorded target.}#1}
\newcommand{\R}{\mathbb R}
\newcommand{\C}{\mathbb C}
\newcommand{\N}{\mathbb N}
\newcommand{\Z}{\mathbb Z}
\newcommand{\Q}{\mathbb Q}
\newcommand{\F}{\mathbb F}
\newcommand{\eps}{\varepsilon}
\begin{document}
% UnsolvedMath ID 30000191; source catalogue code OWR-790-001
\setcounter{section}{30000190}
\section{Systematic Multiscale Design of Hierarchical Materials}\label{30000191}
{\small UnsolvedMath ID 30000191 · Mathematical Physics}
\subsection{Definitions and mathematical statement}

\paragraph{Shared notation.}$\mathbb N=\{1,2,\ldots\}$; the standard number fields have their usual meanings. All additional conventions are stated locally.


\paragraph{Definitions and conventions.} A hierarchical material has structural features at several characteristic length scales $\ell_1<\cdots<\ell_m$, rather than a single microstructure scale. Effective properties, such as stiffness, strength, and flaw tolerance, are macroscopic responses derived from the constitutive and geometric behavior at finer scales. In the report, flaw tolerance means that a pre-existing crack does not propagate before the surrounding material approaches its limiting strength.

\paragraph{Source scope and qualifications.} This is a four-part design research program. The source does not specify a universal class of admissible microstructures, a property map, an optimization norm, or an algorithmic complexity target. These must be chosen for any subsequent precise model; adding them here would invent a narrower problem.

\paragraph{Catalogue mathematical target.}
\paragraph{Catalogue wording retained for traceability.}
How can a hierarchical material be designed systematically to achieve a particular set of properties? How can different levels of structural hierarchy be bridged? What criteria should determine the characteristic length scales at all hierarchical levels? What is the theoretical basis for bottom-up materials design?

\paragraph{Saved literature context.} The supplied catalogue's assessment dated 2026-08-21 records: Hierarchical and flaw-tolerant material models provide design principles and examples, but the general systematic multiscale design, scale-selection criterion, and theoretical bottom-up basis asked here remain open research questions. The assessment is retained as dated provenance; it is not a new comprehensive literature or open-status certification.
\subsection{Short English statement}
Find a systematic theory for designing material structure across many scales to achieve specified mechanical properties.
\subsection{Research attempt}
\paragraph{Model, date and investigated scope.}
OpenAI GPT-6 Astra (configured model \texttt{gpt-6-astra}), ultra reasoning; generated 11 October 2026. Exact series-bar hierarchy calculation demonstrates why stiffness alone cannot determine hierarchy length scales.

\paragraph{Sources and literature check.}
The complete supplied definition, including source qualifications and the dated literature assessment, was reviewed. The calculation below is an elementary restricted attempt motivated by [S2]; it does not claim novelty or a fresh literature-status determination.

\paragraph{Context and investigated scope.} The complete supplied definition, including source qualifications and the dated literature assessment, was reviewed. The calculation below is an elementary restricted attempt motivated by [S2]; it does not claim novelty or a fresh literature-status determination.

\paragraph{A scale-selection obstruction in an exact model.} Consider a one-dimensional elastic bar with positive modulus $E(x)$, length $L$, and imposed end displacement $u(L)-u(0)=eL$. Its energy is $\frac12\int_0^L E(x)|u'(x)|^2dx$. Assume $0<S=\int_0^LE(x)^{-1}dx<\infty$. Cauchy--Schwarz gives $(eL)^2\leq S\int E|u'|^2$, with equality precisely when $u'(x)=eL/(SE(x))$. Therefore
\[
 \min\tfrac12\int_0^L E|u'|^2=\tfrac12LE_{\rm eff}e^2,
 \qquad E_{\rm eff}^{-1}=L^{-1}\int_0^L E^{-1}.
\]
For two constituents in fractions $\theta,1-\theta$, the inverse modulus is $\theta/E_1+(1-\theta)/E_2$, independent of their ordering and absolute layer thicknesses. Applying the same identity inside each hierarchical block and then between blocks simply averages the leaf compliances with their total volume fractions.

\paragraph{Consequences for the proposed design route.} This gives a rigorous bottom-up bridge for the selected elastic response but also proves that stiffness optimization alone cannot select the hierarchy lengths: any rescaling of layers preserving fractions has exactly the same energy. Order the distinct constituent moduli as $0<E_1<E_2$. For a target $E_*\in[E_1,E_2]$, solving $\theta=(E_*^{-1}-E_2^{-1})/(E_1^{-1}-E_2^{-1})$ fixes a fraction, not a length scale. Targets outside that interval are impossible in this series-bar model.

\paragraph{First missing ingredient.} Fracture energy, strength, geometric constraints, or another term with an intrinsic length must enter before the source's flaw-tolerance and scale-selection questions can be answered. This example does not model cracks and does not turn its exact scalar homogenization into a general design theory for multidimensional hierarchical materials. All four research-program questions remain open in their original breadth.

\paragraph{Verification and adversarial audit.}
Checked equality in the energy lower bound, admissible target-modulus interval, identical constituents, and invariance under layer rescaling. Model self-audit only; no independent human proof checking. Prompts A and B were applied to this canonical record and its qualified source scope. Checked equality in the energy lower bound, admissible target-modulus interval, identical constituents, and invariance under layer rescaling. Retained unresolved status and zero original-target progress; A justified fracture/strength model carrying intrinsic length information and a multidimensional design theory are missing.

Coordinator checked weighted Cauchy-Schwarz, the equality displacement, finite-compliance hypothesis, phase fraction and scale invariance in the one-dimensional model. This additional check was model-assisted, not independent human verification.

\paragraph{Tools, process and contributions.}
Python 3 was used for file preparation and structural checks; no numerical experiment or theorem prover was used. Contribution prompt A was followed by prompt B. GPT-6 Astra with ultra reasoning drafted the research and performed its stated self-audit; any additional coordinator contribution is identified in the verification and audit. A GPT-6-based Codex coordinator assembled the manuscript and checked repository integration; its exact variant and reasoning configuration were not exposed. No human mathematical contribution, independent human proof check or formal proof certification is supplied. Complete derivations and the consulted source references are included here; private reasoning transcripts are not published.

\paragraph{Final status and first remaining gap.}
\textbf{UNRESOLVED}. A justified fracture/strength model carrying intrinsic length information and a multidimensional design theory are missing.

\section{Completion Estimate}
\noindent\textbf{COMPLETION: 0\%}
The calculation is an auxiliary or restricted attempt and establishes no newly verified portion of the full source target. A justified fracture/strength model carrying intrinsic length information and a multidimensional design theory are missing.

\subsection{Sources}
\begin{itemize}
\item[\hypertarget{source:30000191:1}{S1}] ULAM.AI and the UnsolvedMath Contributors, UnsolvedMath, record 30000191 (OWR-790-001). Catalogue curation under CC BY 4.0; definitions and English exposition added locally.
\url{https://huggingface.co/datasets/ulamai/UnsolvedMath}.
\item[\hypertarget{source:30000191:2}{S2}] Mini-Workshop: Mathematical Methods and Models of Continuum Biomechanics (2005), Oberwolfach Reports, DOI 10.4171/owr/2005/09, pp. 503--528; official report source and mathematical context.
\url{https://doi.org/10.4171/owr/2005/09}.
\item[\hypertarget{source:30000191:3}{S3}] Official EMS Press PDF of the same Oberwolfach report; PDF page 6, printed page 508.
\url{https://ems.press/content/serial-article-files/45983}.
\item[\hypertarget{source:30000191:4}{S4}] Gao et al., materials become insensitive to flaws at nanoscale; reference in the saved dated catalogue literature review.
\url{https://doi.org/10.1073/pnas.0937635100}.
\end{itemize}
\label{end:30000191}
\end{document}
